% ==============================================================================
% 08_criteria5_staff.tex — Criteria 5: Academic Staff
% Score: 2
% ==============================================================================

\criterion{5}{บุคลากรสายวิชาการ}{Academic Staff}

% ------------------------------------------------------------------------------
\criterionsub{5.1 The programme to show that academic staff planning (including succession, promotion, re-deployment, termination, and retirement plans) is carried out to ensure that the quality and quantity of the academic staff fulfil the needs for education, research, and service.}

หลักสูตรวิศวกรรมศาสตรมหาบัณฑิต สาขาวิชาวิศวกรรมคอมพิวเตอร์และปัญญาประดิษฐ์ (หลักสูตรใหม่ พ.ศ.~2568) ปีการศึกษา 2568 กำหนดให้มี\textbf{อาจารย์ผู้รับผิดชอบหลักสูตรจำนวน 3 คน} ได้แก่ ผศ.ดร.พิศิษฐ์ นาคใจ ผศ.ดร.กาญจนา ดาวเด่น และ ผศ.ดร.พัชรี มณีรัตน์ ทุกท่านมีคุณวุฒิปริญญาเอก (ปร.ด.) และตำแหน่งทางวิชาการระดับผู้ช่วยศาสตราจารย์ เป็นไปตามเกณฑ์มาตรฐานบัณฑิตศึกษา โดยผ่านการสรรหาและการกำหนดอัตรากำลังภายใต้กรอบของมหาวิทยาลัยและคณะเทคโนโลยีอุตสาหกรรม นอกจากนี้หลักสูตรยังประสาน\textbf{อาจารย์ผู้สอนและอาจารย์ที่เกี่ยวข้อง}จากภายในมหาวิทยาลัยตามความจำเป็นของรายวิชาและงานวิจัย ตามรายชื่อในส่วนข้อมูลประกอบของ SAR

มหาวิทยาลัยดำเนินงานด้านการบริหารทรัพยากรบุคคลภายใต้แผนบริหารทรัพยากรบุคคลประจำปี พ.ศ.~2565--2569 ครอบคลุมการสืบทอดตำแหน่ง การเลื่อนขั้น การโยกย้าย การเลิกจ้าง และแผนเกษียณอายุ คณะกรรมการหลักสูตรและคณะกรรมการบัณฑิตศึกษาร่วมกับงานบุคคล กำหนดแผนสืบทอดบทบาทหัวหน้าหลักสูตร อาจารย์ที่ปรึกษารายวิชาหลัก และการเสริมอาจารย์ผู้สอนจากภายในมหาวิทยาลัยตามความจำเป็นของรายวิชาและงานวิจัย เพื่อลดความเสี่ยงจากการเกษียณหรือการโยกย้ายของบุคลากร รายละเอียดประวัติและผลงานของอาจารย์ผู้รับผิดชอบหลักสูตรและอาจารย์ผู้สอนที่เกี่ยวข้องสรุปไว้ในส่วนข้อมูลประกอบของ SAR

การดำเนินการเพิ่มเติมที่สอดคล้องกับแนวทางของหลักสูตรอื่นในมหาวิทยาลัย (เช่น การทบทวนแผนพัฒนาบุคลากรรายบุคคลและการสื่อสารนโยบายการเข้าสู่ตำแหน่งทางวิชาการ) หลักสูตรจะจัดทำและแนบเป็นหลักฐาน AUNQA เมื่อเริ่มมีการประชุมหลักสูตรรายภาคการศึกษาอย่างต่อเนื่อง

ทั้งนี้ \textbf{ผู้รับผิดชอบหลัก} คือกองบริหารงานบุคคล มหาวิทยาลัยราชภัฏอุตรดิตถ์ และคณะเทคโนโลยีอุตสาหกรรม (งานบัณฑิตศึกษา/คณะกรรมการหลักสูตร) โดยใช้ \textbf{หลักฐานประกอบ} ได้แก่ AUNQA-5-1 แผนบริหารทรัพยากรบุคคล พ.ศ.~2565--2569 และตารางอาจารย์ผู้รับผิดชอบหลักสูตรในส่วนข้อมูลประกอบของ SAR

% ------------------------------------------------------------------------------
\criterionsub{5.2 The programme to show that staff workload is measured and monitored to improve the quality of education, research, and service.}

หลักสูตรให้ความสำคัญกับการวัดและติดตามภาระงานของอาจารย์ เพื่อคุณภาพการจัดการเรียนการสอน การวิจัย และการบริการวิชาการ โดยอ้างอิงประกาศมหาวิทยาลัยเรื่องแนวทางการคำนวณค่าภาระงานเทียบเท่าเต็มเวลาของอาจารย์ (FTE) และค่าภาระการเรียนของผู้เรียน ตามที่งานมาตรฐานการศึกษาและประกันคุณภาพการศึกษาเผยแพร่ (เช่น \href{https://academic.uru.ac.th/document/06912.pdf}{academic.uru.ac.th/document/06912.pdf}) ร่วมกับกองบริหารงานบุคคล โดยพิจารณาทั้งภาระสอน วิจัย บริการวิชาการ และทำนุบำรุงศิลปวัฒนธรรม คณะกรรมการหลักสูตรทบทวนความสมดุลของภาระงานสอนรายวิชาในกลุ่มวิชาหลักและกลุ่มวิชาเลือกทุกภาคการศึกษา เพื่อไม่ให้ภาระงานสอนรวมของอาจารย์แต่ละท่านเกินเกณฑ์ที่มหาวิทยาลัยกำหนด และให้มีเวลาสำหรับการที่ปรึกษาวิทยานิพนธ์/สารนิพนธ์

\begin{table}[H]
  \centering
  \caption{จำนวนอาจารย์ผู้รับผิดชอบหลักสูตรและการติดตามภาระงาน (FTE) ของหลักสูตร CE\&AI}
  \begin{tabularx}{\linewidth}{|l|c|Y|Y|}
    \hline
    \rowcolor{gray!15}
    \textbf{ปีการศึกษา} & \textbf{จำนวนอาจารย์ผู้รับผิดชอบ} & \textbf{FTE ระดับบุคลากรประจำมหาวิทยาลัย} & \textbf{การติดตามภาระงานในหลักสูตร} \\
    \hline
    2568 & 3 คน & 3.0 (อาจารย์ผู้รับผิดชอบหลักสูตร ประจำเต็มเวลา) & สรุปสัดส่วนภาระงานสอน/ที่ปรึกษาเฉพาะ CE\&AI หลังมีตารางสอนภาคแรก \\
    \hline
  \end{tabularx}
\end{table}

\marknote{ค่า FTE 3.0 หมายถึงจำนวนเทียบเท่าการจ้างเต็มเวลาของอาจารย์ผู้รับผิดชอบหลักสูตร 3 ท่านในระดับมหาวิทยาลัย; ตัวเลขแยกตามภาระงานสอนในหลักสูตร CE\&AI เทียบกับหลักสูตรอื่นจะบันทึกหลังเริ่มเปิดสอนและมีข้อมูลตารางสอนจริง}

% ------------------------------------------------------------------------------
\criterionsub{5.3 The programme to show that the competences of the academic staff are determined, evaluated, and communicated.}

มหาวิทยาลัยโดยกองบริหารงานบุคคลจัดทำประกาศมหาวิทยาลัยราชภัฏอุตรดิตถ์ เรื่อง นโยบายและแนวทางการพัฒนาคุณภาพอาจารย์ ซึ่งกำหนดสมรรถนะที่อาจารย์พึงมี (Core competence + Functional competence) และมีการประเมินผ่านระบบ TPS (Teacher Performance System) และการประเมินสมรรถนะรายปี ดังนี้
\begin{enumerate}[leftmargin=2.2em,itemsep=2pt,topsep=2pt]
  \item \textbf{การต่อสัญญาจ้าง/การประเมินสมรรถนะ} --- ดำเนินการตามประกาศมหาวิทยาลัยที่เกี่ยวข้อง และเผยแพร่ผ่านเว็บไซต์สภามหาวิทยาลัย/งานบุคคล (ตัวอย่างช่องทางสารสนเทศ: \href{https://personnel.uru.ac.th/}{personnel.uru.ac.th})
  \item \textbf{การเลื่อนขั้นเงินเดือนและผลการปฏิบัติราชการ} --- มหาวิทยาลัยกำหนดรอบการประเมินเป็นประจำ (โดยทั่วไปแบ่งเป็นรอบที่ 1: 1 ตุลาคม--31 มีนาคม และรอบที่ 2: 1 เมษายน--30 กันยายน) ผู้ประเมินสายวิชาการรวมถึงประธานหลักสูตร ผู้ช่วยคณบดี รองคณบดี และคณบดี ผลการประเมินมีการสื่อสารผ่านระบบ HRMS และช่องทางของคณะตามแนวปฏิบัติของมหาวิทยาลัย
\end{enumerate}
ผลการประเมินและข้อเสนอแนะเพื่อการพัฒนามีการสื่อสารถึงอาจารย์ผ่านผู้บังคับบัญชาและระบบสารสนเทศของมหาวิทยาลัย

หลักฐานประกอบเกณฑ์ย่อยนี้ประกอบด้วย AUNQA-5-2 ประกาศมหาวิทยาลัยฯ เรื่อง นโยบายและแนวทางการพัฒนาคุณภาพอาจารย์ และบันทึกผลการประเมินสมรรถนะ/TPS ของอาจารย์ผู้รับผิดชอบหลักสูตร (เมื่อรวบรวมครบตามรอบการประเมิน)

% ------------------------------------------------------------------------------
\criterionsub{5.4 The programme to show that the duties allocated to the academic staff are appropriate to qualifications, experience, and aptitude.}

อาจารย์มหาวิทยาลัยมีภาระหน้าที่หลักด้านการสอน การวิจัย การบริการวิชาการ และทำนุบำรุงศิลปวัฒนธรรม หลักสูตร CE\&AI จัดสรรภาระงานสอน การที่ปรึกษาวิทยานิพนธ์/สารนิพนธ์ และงานวิจัยตามคุณวุฒิ ประสบการณ์ และความถนัดของ\textbf{อาจารย์ผู้รับผิดชอบหลักสูตร 3 คน} โดยสรุปบทบาทหลักในหลักสูตรเป็นตารางที่ \ref{tbl:c5_roles} และรายละเอียดความเชี่ยวชาญตามคุณวุฒิเป็นย่อหน้าถัดไป

\begin{table}[H]
  \centering
  \caption{บทบาทหลักของอาจารย์ผู้รับผิดชอบหลักสูตร CE\&AI จำนวน 3 คน (สอดคล้องการจัดสรรงานตามความถนัด)}
  \label{tbl:c5_roles}
  \begingroup
  \small
  \setlength{\tabcolsep}{4pt}
  \renewcommand{\arraystretch}{1.2}
  \begin{tabularx}{\linewidth}{|>{\raggedright\arraybackslash}p{3.8cm}|Y|}
    \hline
    \rowcolor{gray!15}
    \textbf{อาจารย์ผู้รับผิดชอบหลักสูตร} & \textbf{บทบาท/ความรับผิดชอบหลักที่เกี่ยวข้องกับหลักสูตร} \\
    \hline
    ผศ.ดร.พิศิษฐ์ นาคใจ & กลุ่มวิชาปัญญาประดิษฐ์/การเรียนรู้ของเครื่อง การประมวลผลภาพและสัญญาณ \\
    \hline
    ผศ.ดร.กาญจนา ดาวเด่น & ระบบฝังตัว ระบบอัตโนมัติ และการประยุกต์ใช้ในภาคสนาม \\
    \hline
    ผศ.ดร.พัชรี มณีรัตน์ & สถิติและการวิเคราะห์ข้อมูล ระเบียบวิธีวิจัยเชิงปริมาณ \\
    \hline
  \end{tabularx}
  \endgroup
\end{table}

\noindent\textbf{คุณวุฒิและความเชี่ยวชาญที่ใช้จัดสรรงาน}
ได้แก่ ผศ.ดร.พิศิษฐ์ นาคใจ (ปร.ด. วิศวกรรมคอมพิวเตอร์) ซึ่งมีความเชี่ยวชาญด้านปัญญาประดิษฐ์และการเรียนรู้ของเครื่อง รวมถึงการประมวลผลภาพและสัญญาณ; ผศ.ดร.กาญจนา ดาวเด่น (ปร.ด. วิศวกรรมคอมพิวเตอร์) ซึ่งเชี่ยวชาญด้านระบบฝังตัว ระบบอัตโนมัติ และการประยุกต์ใช้ในภาคสนาม; และ ผศ.ดร.พัชรี มณีรัตน์ (ปร.ด. สถิติประยุกต์) ซึ่งเชี่ยวชาญด้านสถิติและการวิเคราะห์ข้อมูล รวมทั้งระเบียบวิธีวิจัยเชิงปริมาณ

% ------------------------------------------------------------------------------
\criterionsub{5.5 The programme to show that promotion of the academic staff is based on a merit system which accounts for teaching, research, and service.}

มหาวิทยาลัยใช้ระบบเลื่อนตำแหน่งและการประเมินผลการปฏิบัติราชการตามเกณฑ์ ก.พ.อ. โดยยึดผลงานด้านการสอน การวิจัย และการบริการวิชาการ (รวมการทำนุบำรุงศิลปวัฒนธรรมและการพัฒนาตนเองตามแบบประเมินของมหาวิทยาลัย) มีกรรมการประเมินผลงานทางวิชาการประจำมหาวิทยาลัย และมีการประกาศแนวทางการเข้าสู่ตำแหน่งทางวิชาการให้อาจารย์ทราบอย่างสม่ำเสมอ อาจารย์ผู้รับผิดชอบหลักสูตร CE\&AI จำนวน 3 คน ทุกท่านเป็นผู้ช่วยศาสตราจารย์และมีคุณวุฒิปริญญาเอก ซึ่งได้รับการแต่งตั้ง/เลื่อนตำแหน่งผ่านระบบ merit ตามหลักเกณฑ์ดังกล่าว

% ------------------------------------------------------------------------------
\criterionsub{5.6 The programme to show that the rights and privileges, benefits, roles and relationships, and accountability of the academic staff, taking into account professional ethics and their academic freedom, are well defined and understood.}

สิทธิและสิทธิพิเศษ บทบาทและความสัมพันธ์ และความรับผิดชอบของอาจารย์ ระบุไว้ใน
ข้อบังคับมหาวิทยาลัยราชภัฏอุตรดิตถ์ ว่าด้วยการบริหารงานบุคคลของพนักงานมหาวิทยาลัย ประมวลจริยธรรมของบุคลากรมหาวิทยาลัย และคู่มืออาจารย์/พนักงานมหาวิทยาลัย โดยทั้งหมดคำนึงถึงจริยธรรมทางวิชาชีพและเสรีภาพทางวิชาการ (academic freedom)
ทั้งคำนึงถึงจริยธรรมทางวิชาชีพและเสรีภาพทางวิชาการ (academic freedom)

มหาวิทยาลัยมีการกำหนดภาระงานขั้นต่ำของอาจารย์ประจำและช่องทางสื่อสารสิทธิประโยชน์ตามประกาศที่เกี่ยวข้อง (อ้างอิงเช่น \href{https://academic.uru.ac.th/document/06912.pdf}{ประกาศภาระงานขั้นต่ำของคณาจารย์ประจำ} และประกาศสวัสดิการที่เผยแพร่ผ่านเว็บไซต์สภามหาวิทยาลัย) หลักสูตรสื่อสารบทบาทอาจารย์ที่ปรึกษาและจรรยาบรรณทางวิชาชีพในระดับหลักสูตรและคณะให้อาจารย์และนักศึกษาทราบเมื่อเริ่มเปิดสอน

% ------------------------------------------------------------------------------
\criterionsub{5.7 The programme to show that the training and developmental needs of the academic staff are systematically identified, and that appropriate training and development activities are implemented to fulfil the identified needs.}

ในทุกปีการศึกษา หลักสูตรจะดำเนินการทบทวนและสำรวจความต้องการพัฒนาอาจารย์ผ่านแบบสำรวจของมหาวิทยาลัยและการประชุมคณะกรรมการหลักสูตร แจ้งแหล่งงบประมาณสนับสนุนจากคณะ/มหาวิทยาลัย และประสานให้อาจารย์เข้าร่วมกิจกรรมที่จัดเป็นประจำ ตัวอย่างกิจกรรมที่สอดคล้องกับลักษณะหลักสูตร CE\&AI ได้แก่
การอบรม AUN-QA และการเป็นผู้ประเมินภายนอก (Assessor) ตามแผนพัฒนาคุณภาพการศึกษา การอบรมการจัดการเรียนการสอนแบบ OBE และ Constructive Alignment รวมถึงการอบรม Active Learning ระเบียบวิธีวิจัย และการใช้เครื่องมือ AI เพื่อยกระดับประสิทธิผลการจัดการเรียนการสอนและการวิจัย

% ------------------------------------------------------------------------------
\criterionsub{5.8 The programme to show that performance management including reward and recognition is implemented to assess academic staff teaching and research quality.}

มหาวิทยาลัยมีระบบ Performance Management ผ่าน TPS (Teacher Performance System) ประเมินคุณภาพการสอนและการวิจัยรายปี รวมถึงการติดตามแผนพัฒนาอาจารย์รายบุคคล และมีรางวัลอาจารย์ดีเด่น ผลงานวิจัยดีเด่น เพื่อสร้างแรงจูงใจในการพัฒนาผลงานทั้งสามด้าน

\begin{selfassessmentbox}{2}{อาจารย์ผู้รับผิดชอบหลักสูตร 3 คน (ผศ.ดร.พิศิษฐ์ นาคใจ ผศ.ดร.กาญจนา ดาวเด่น ผศ.ดร.พัชรี มณีรัตน์) คุณวุฒิปริญญาเอกทั้งหมดและตำแหน่งทางวิชาการระดับผู้ช่วยศาสตราจารย์ มีระบบบริหารบุคคล การประเมินสมรรถนะ และการพัฒนาอาจารย์ตามนโยบายมหาวิทยาลัยครบถ้วน มีการประสานอาจารย์ผู้สอนร่วมตามรายวิชาในส่วนข้อมูลประกอบ ส่วนข้อมูลภาระงานสอนเฉพาะหลักสูตร CE\&AI แยกจากหลักสูตรอื่นจะสมบูรณ์หลังมีตารางสอนและข้อมูลการปฏิบัติงานรายภาคการศึกษาจริง}
\end{selfassessmentbox}

% ------------------------------------------------------------------------------
\subsection*{รายการหลักฐาน}

หลักฐานประกอบ Criterion~5 ประกอบด้วย AUNQA-5-1 แผนบริหารทรัพยากรบุคคล มหาวิทยาลัยราชภัฏอุตรดิตถ์ พ.ศ.~2565--2569, AUNQA-5-2 ประกาศมหาวิทยาลัยฯ เรื่อง นโยบายและแนวทางการพัฒนาคุณภาพอาจารย์, AUNQA-5-3 ข้อบังคับมหาวิทยาลัยฯ ว่าด้วยการบริหารงานบุคคลของพนักงานมหาวิทยาลัย และ AUNQA-5-4 รายงานการพัฒนาอาจารย์/บันทึกการเข้าร่วมอบรมของอาจารย์ผู้รับผิดชอบหลักสูตร
